{\bfseries Algoritmo de Simon con una función de paridad modificada.

Sea $f : \{0,1\}^3 \to \{0,1\}^3$ definida por

$$
    f(x) =
    \begin{cases}
        x,            & \text{si } x \text{ tiene peso par},   \\
        x \oplus 101, & \text{si } x \text{ tiene peso impar}.
    \end{cases}
$$
}

{\bfseries (a) Verifique que $f$ es 2 a 1 y determine su período $s$ tal que $f(x)=f(x\oplus s)$ para todo $x$.}

Evaluamos $f$ aplicando correctamente el peso de Hamming (número de unos):

\begin{center}
    \begin{tabular}{c|c|c||c|c|c}
        $x$   & peso      & $f(x)$                & $x$   & peso      & $f(x)$                \\ \hline
        $000$ & 0 (par)   & $000$                 & $100$ & 1 (impar) & $100\oplus 101 = 001$ \\
        $011$ & 2 (par)   & $011$                 & $101$ & 2 (par)   & $101$                 \\
        $001$ & 1 (impar) & $001\oplus 101 = 100$ & $110$ & 2 (par)   & $110$                 \\
        $010$ & 1 (impar) & $010\oplus 101 = 111$ & $111$ & 3 (impar) & $111\oplus 101 = 010$ \\
    \end{tabular}
\end{center}

Las ocho imágenes son $\{000, 011, 100, 111, 001, 101, 110, 010\} = \{0,1\}^3$, todas distintas. Por lo tanto $f$ es \textbf{inyectiva} (1 a 1), no 2 a 1.

La razón es estructural: XOR con $s = 101$ es una permutación de $\{0,1\}^3$. Para que la función fuese 2 a 1 se necesitaría que $x \oplus s$ cambie de paridad respecto a $x$, lo cual ocurre solo si $s$ tiene peso \emph{impar}; pero $|101|_H = 2$ (par), por lo que $x$ y $x\oplus 101$ siempre comparten la misma paridad y caen en ramas distintas de la función.

El único $s$ que satisface $f(x) = f(x \oplus s)$ para \emph{todo} $x$ es:
$$
    \boxed{s = 000.}
$$

\vspace{0.5cm}

{\bfseries (b) Simule una ejecución del algoritmo de Simon: inicialice $|000\rangle \otimes |000\rangle$, aplique Hadamard a ambos registros, luego el oráculo $U_f$ definido por}

$$
    U_f |x\rangle |y\rangle = |x\rangle |y \oplus f(x)\rangle,
$$

{\bfseries y mida el segundo registro. Suponga que el resultado es $f(000)=000$. Determine el estado colapsado del primer registro, aplique la segunda Hadamard y obtenga los posibles resultados y sus probabilidades. Verifique que $z \cdot s = 0 \pmod 2$.}

\textbf{Inicialización y primera Hadamard.} Partimos de $\ket{000}\ket{000}$. Tras $H^{\otimes 3}$ en el primer registro:
$$
    \ket{\psi_1} = \frac{1}{2\sqrt{2}} \sum_{x \in \{0,1\}^3} \ket{x}\ket{000}.
$$

\textbf{Oráculo.} Aplicamos $U_f$:
$$
    \ket{\psi_2} = \frac{1}{2\sqrt{2}} \sum_{x} \ket{x}\ket{f(x)}.
$$

\textbf{Medición del segundo registro.} Como $f$ es inyectiva, solo $x = 000$ cumple $f(x) = 000$. El primer registro colapsa a:
$$
    \ket{\psi_{\text{col}}} = \ket{000}.
$$

\textbf{Segunda Hadamard.} Aplicamos $H^{\otimes 3}$ a $\ket{000}$:
$$
    H^{\otimes 3}\ket{000} = \frac{1}{2\sqrt{2}} \sum_{z \in \{0,1\}^3} \ket{z}.
$$
Los ocho resultados posibles son $\{000, 001, 010, 011, 100, 101, 110, 111\}$, cada uno con amplitud $\frac{1}{2\sqrt{2}}$ y probabilidad:
$$
    P(z) = \left(\frac{1}{2\sqrt{2}}\right)^{\!2} = \frac{1}{8}
$$

\textbf{Verificación} $z \cdot s = 0 \pmod{2}$ con $s = 000$: para cualquier $z$,
$$
    z \cdot 000 = z_1 \cdot 0 \oplus z_2 \cdot 0 \oplus z_3 \cdot 0 = 0. \quad
$$
Todos los $z$ satisfacen la condición, consistente con que $s = 000$ (función inyectiva).

\vspace{0.5cm}

{\bfseries (c) ¿Cuál es la probabilidad de obtener dos vectores distintos y no nulos en dos ejecuciones que permitan determinar $s$? Considere que la distribución es uniforme en el subespacio $z \cdot s = 0$.}

Con $s = 000$, el subespacio es $S^\perp = \{z : z \cdot 000 = 0\} = \{0,1\}^3$ completo, con $2^3 = 8$ elementos equiprobables ($P = 1/8$ cada uno).

\begin{itemize}
    \item \textbf{Primera ejecución:} probabilidad de obtener $z_1 \neq 000$: $\dfrac{7}{8}$.
    \item \textbf{Segunda ejecución:} dado $z_1 \neq 000$, necesitamos $z_2 \notin \{000, z_1\}$ para que sean independientes. Hay $8 - 2 = 6$ opciones favorables de 8 posibles: $\dfrac{6}{8} = \dfrac{3}{4}$.
\end{itemize}

La probabilidad conjunta es:
$$
    \boxed{P = \frac{7}{8} \times \frac{3}{4} = \frac{21}{32}.}
$$

Obtener dos vectores independientes forma un sistema de $n-1 = 2$ ecuaciones. Al quedar una variable libre, el sistema produce exactamente dos candidatos: la solución trivial $s = 000$ y una solución $s' \neq 000$. Este es el resultado esperado de la fase cuántica del Algoritmo de Simon. La determinación de $s$ se concluye mediante el paso clásico final: evaluar y comparar $f(000)$ con $f(s')$. Como $f(000) \neq f(s')$, se descarta $s'$ y se determina con certeza que $s = 000$.